%%═══════════════════════════════════════════════════════════════════
%% Lecture 16 — Multigrid Methods
%% 77315 — Computational Physics
%%═══════════════════════════════════════════════════════════════════

\section{Lecture 16 — Multigrid Methods}\label{sec:lec16}
\hrule\vspace{1.5em}

%%──────────────────────────────────────────────────────────────────
\subsection*{Overview}
SOR with $\omega_{\rm opt}$ still requires $\bigO{N}$ iterations,
giving $\bigO{N^3}$ total work on an $N\times N$ grid.
\emph{Multigrid} achieves $\bigO{N^2}$ (optimal) by exploiting the
observation that relaxation quickly damps high-frequency errors
but is slow for low-frequency modes — which look high-frequency on
a coarser grid.

%%──────────────────────────────────────────────────────────────────
\subsection{Error, Residual, and Correction}\label{subsec:lec16:err}

Consider the linear system $\mathcal{L}\mathbf{u} = \mathbf{f}$
arising from discretising a linear elliptic PDE.  After a few
relaxation sweeps, the current approximation $\tilde{\mathbf{u}}$
has error $\mathbf{e} = \mathbf{u}^* - \tilde{\mathbf{u}}$
and residual
\begin{equation}
  \mathbf{r} = \mathbf{f} - \mathcal{L}\tilde{\mathbf{u}}.
  \label{eq:lec16:residual}
\end{equation}
The error satisfies the \emph{residual equation}:
\begin{equation}
  \mathcal{L}\mathbf{e} = \mathbf{r}.
  \label{eq:lec16:correction}
\end{equation}
Solve~\eqref{eq:lec16:correction} on a \emph{coarser} grid to get an
approximate correction, then add it back to $\tilde{\mathbf{u}}$.

%%──────────────────────────────────────────────────────────────────
\subsection{Restriction and Prolongation Operators}\label{subsec:lec16:ops}

\subsubsection*{Restriction $\mathcal{R}$: fine $\to$ coarse}
Transfer the residual from the fine grid ($h$) to a coarse grid ($2h$).
Full-weighting restriction in 2D:
\begin{equation}
  r^{2h}_{I,J} = \frac{1}{16}\left[
    4r^h_{2I,2J}
    + 2(r^h_{2I\pm1,2J}+r^h_{2I,2J\pm1})
    + r^h_{2I\pm1,2J\pm1}
  \right].
  \label{eq:lec16:restrict}
\end{equation}

\subsubsection*{Prolongation $\mathcal{P}$: coarse $\to$ fine}
Interpolate (prolong) the coarse-grid correction back to the fine grid.
Bilinear interpolation assigns weights to the four surrounding coarse
nodes based on the fractional distance.  For a fine-grid point at the
midpoint between two coarse nodes, the weight is $1/2$; at the
midpoint of a coarse cell, the weight is $1/4$.

%%──────────────────────────────────────────────────────────────────
\subsection{V-Cycle Algorithm}\label{subsec:lec16:vcycle}

\begin{enumerate}
  \item \textbf{Pre-smoothing}: perform $\nu_1$ relaxation sweeps on
        the fine grid.
  \item Compute residual $\mathbf{r} = \mathbf{f} - \mathcal{L}\tilde{\mathbf{u}}$.
  \item Restrict: $\mathbf{r}^{2h} = \mathcal{R}\,\mathbf{r}$.
  \item \textbf{Coarse-grid solve}: recursively apply the V-cycle on
        the coarser grid to solve $\mathcal{L}^{2h}\mathbf{e}^{2h}=\mathbf{r}^{2h}$.
        On the coarsest grid, solve exactly.
  \item Prolongate and correct: $\tilde{\mathbf{u}} \leftarrow
        \tilde{\mathbf{u}} + \mathcal{P}\,\mathbf{e}^{2h}$.
  \item \textbf{Post-smoothing}: perform $\nu_2$ relaxation sweeps.
\end{enumerate}

The V-cycle makes one recursive call per level; see figure below.

\subsubsection*{W-cycle}
The W-cycle makes \emph{two} recursive calls per level,
giving $\bigO{N^2 \log N}$ work but more robust error reduction.
For most problems the V-cycle with $\nu_1=\nu_2=1$ is sufficient.

%%──────────────────────────────────────────────────────────────────
\subsection{Convergence}\label{subsec:lec16:convergence}

Each V-cycle reduces the error by a constant factor $\rho < 1$
independent of the grid size $N$.  Hence the total work to reach
tolerance $\eps$ is $\bigO{N^2 \log(1/\eps)}$ — essentially
optimal (compare $\bigO{N^4}$ for Jacobi).

\nt{Key insight: relaxation damps \emph{high-frequency} errors fast.
After restriction, the remaining low-frequency error looks
high-frequency on the coarser grid and is quickly damped there.
This is the core of why multigrid works.}
